\chapter{اللقاء الثالث والعشرون: ختام أحكام الصيام والاعتكاف وبداية آيات القتال}

\section{مقدمة}
السلام عليكم ورحمة الله وبركاته. الحمد لله رب العالمين، وأشهد أن لا إله إلا الله وحده لا شريك له، وأشهد أن محمداً عبده ورسوله. اللهم صل على محمد وعلى آل محمد كما صليت على آل إبراهيم إنك حميد مجيد، اللهم بارك على محمد وعلى آل محمد كما باركت على آل إبراهيم إنك حميد مجيد.

استأنف الشيخ هذا اللقاء من قوله تعالى: \begin{ayah}يُرِيدُ اللَّهُ بِكُمُ الْيُسْرَ وَلَا يُرِيدُ بِكُمُ الْعُسْرَ\end{ayah}، بعد تمام الكلام على رخص الصيام، ثم مضى في آيات الدعاء، والاعتكاف، والأهلة، والقتال.

\separator

\section{تأويل (يريد الله بكم اليسر ولا يريد بكم العسر)}
\isnad{قال أبو جعفر:} يريد الله بالمؤمنين التخفيف والتيسير فيما شرعه لهم من الفطر في المرض والسفر مع وجوب القضاء بعد ذلك، ولا يريد بهم المشقة والتكليف بما يشق عليهم في تلك الأحوال.

وساق \rawi{الطبري} آثاراً عن \rawi{ابن عباس} و\rawi{مجاهد} و\rawi{قتادة} وغيرهم أن المراد باليسر هنا: الرخصة في الفطر في السفر والمرض، وبالعسر: إلزام الصيام في تلك الحال.

\begin{fawaid}[إرادة العبد لنفسه ينبغي أن تكون تابعة لمراد الله به]
من أجلّ المعاني في هذا الموضع قول \rawi{قتادة}: ``فأريدوا لأنفسكم ما أراد الله بكم''. فكلما تعلّم العبد ما أراده الله به في الفعل والترك، ثم حمل نفسه عليه، كان ذلك من أعظم أبواب الرشد وتزكية النفس.
\end{fawaid}

\separator

\section{تأويل (ولتكملوا العدة)}
\isnad{قال أبو جعفر:} معنى الآية: لتكملوا عدة ما أفطرتم من رمضان في السفر والمرض بقضائه من أيام أخر بعد البرء والإقامة.

وذكر \rawi{الطبري} في هذا الموضع خلافاً لغوياً في تعلّق اللام والواو، ونقل فيه عن \rawi{الأخفش} و\rawi{الفراء}، ثم رجح ما يراه أوضح في العربية: أن اللام هنا متعلقة بفعل مقدر بعدها، لا معطوفةً على ما قبلها على الوجه الذي ذكره بعضهم.

\begin{fawaid}[الترجيح اللغوي عند الطبري ليس نقلاً محضاً]
لا يكتفي الطبري في المسائل اللغوية بحكاية أقوال النحاة، بل يوازن بينها في ضوء تركيب الآية نفسها، ويدفع ما يراه أضعف من جهة الصناعة أو السياق. وهذا يبين أن عنايته بالعربية عناية استدلال وترجيح، لا مجرد جمعٍ للأوجه.
\end{fawaid}

\section{تأويل (ولتكبروا الله على ما هداكم)}
فسر \rawi{الطبري} ذلك بالتكبير يوم الفطر، وذكر آثاراً في ابتداء التكبير من رؤية هلال شوال إلى خروج الإمام للصلاة وفراغه منها.

\section{تأويل (ولعلكم تشكرون)}
ذكر أن المراد: لتشكروا الله على نعمة الهداية والتوفيق والتيسير، وأن ``لعل'' هنا أقرب ما تكون إلى معنى ``كي''.

\separator

\section{تأويل (وإذا سألك عبادي عني فإني قريب)}
\isnad{قال أبو جعفر:} إذا سألك عبادي عني أين أنا أو متى أدعى أو كيف أدعى، فإني قريب منهم، أسمع دعاءهم وأجيب من دعاني إذا دعاني.

وجمع \rawi{الطبري} أسباباً متعددة للنزول في الآية:
\begin{itemize}
    \item سؤال بعضهم: أقريب ربنا فنناجيه أم بعيد فنناديه؟
    \item سؤالهم: أي ساعة ندعو الله فيها؟
    \item سؤالهم: كيف ندعوه؟
\end{itemize}

ورأى أن هذه المعاني متقاربة، وأن الآية جواب جامع لها.

ثم نبه الشيخ إلى فائدة دقيقة في صنيع \rawi{الطبري}: أنه لا يجعل تعدد المرويات موجباً للاضطراب دائماً، بل قد يراها متعاضدة إذا أمكن حملها على معنى كلي واحد، كما في هذه الآية التي دارت أسئلة الناس فيها حول قرب الله، ووقت دعائه، وكيفية سؤاله.

\begin{fawaid}[الأولى في باب الدعاء سؤال الإخلاص لا مجرد سؤال الوقت]
نبّه الشيخ هنا إلى معنى نفيس: أن كثيراً من الناس ينشغل بالسؤال عن وقت الإجابة، والأهم من ذلك أن يسأل عن صدق الدعاء وإخلاصه وحضور قلبه؛ لأن الإشكال الأعظم في كثير من الأحوال ليس في الوقت، بل في حقيقة التوجه إلى الله.
\end{fawaid}

\begin{fawaid}[تعدد أسباب النزول لا يعني تعدد المعنى بالضرورة]
إذا وردت في الآية أخبار متعددة في سبب السؤال أو المناسبة، وكان يمكن جمعها تحت جواب قرآني واحد، فالأولى حملها على التوافق لا على التضاد. وهذا من حسن فقه الطبري في التعامل مع الروايات التفسيرية.
\end{fawaid}

\separator

\section{تأويل (أحل لكم ليلة الصيام الرفث إلى نسائكم)}
\isnad{قال أبو جعفر:} أطلق الله لكم في ليلة الصيام الجماع وسائر ما يتعلق به بعد أن كان قد حُظر على بعض من قبلهم أو على بعض أحوال المسلمين أول الإسلام.

\subsection{تأويل (هن لباس لكم وأنتم لباس لهن)}
ذكر \rawi{الطبري} في معنى اللباس وجهين:
\begin{itemize}
    \item أنه كناية عن شدة الملابسة والمخالطة والاتصال.
    \item وأن كل واحد من الزوجين ستر وسكن للآخر.
\end{itemize}

\subsection{تأويل (علم الله أنكم كنتم تختانون أنفسكم)}
بين \rawi{الطبري} أن الله علم ما كان يقع من بعض المؤمنين من مخالفة هذا الحكم قبل تخفيفه، فتاب عليهم وعفا عنهم، وأباح لهم ما كانوا يتكلفون تركه.

\section{تأويل (وكلوا واشربوا حتى يتبين لكم الخيط الأبيض من الخيط الأسود من الفجر)}
استوعب \rawi{الطبري} الأقوال في تفسير الخيط الأبيض والخيط الأسود، وذكر آثار \rawi{عدي بن حاتم} وما يتعلق بفهم الآية، حتى قرر أن المراد بهما بياض الفجر وسواد الليل.

\section{تأويل (ثم أتموا الصيام إلى الليل)}
فسر \rawi{الطبري} الآية بأن الواجب هو الإمساك من تبين الفجر إلى دخول الليل، ثم يباح بعد ذلك الأكل والشرب وسائر ما كان ممنوعاً في زمن الصيام.

وفي أثناء هذا الموضع ذكر \rawi{الطبري} الأخبار المتظاهرة في النهي عن الوصال، وأن غاية ما رخص فيه بعضهم هو الوصال إلى السحر لا أكثر.

\begin{fawaid}[استعمال الطبري للتواتر أوسع من اصطلاح المتكلمين]
استعمل الطبري هنا لفظ ``الأخبار المتواترة'' في الأخبار الكثيرة الثابتة المشهورة عن النبي ﷺ، من غير التقييد الاصطلاحي المتأخر الذي شدد فيه بعض المتكلمين. وهذا يبين أن استعمال السلف للمصطلحات أوسع وأقرب إلى المعنى اللغوي والعملي.
\end{fawaid}

\separator

\section{تأويل (ولا تباشروهن وأنتم عاكفون في المساجد)}
\isnad{قال أبو جعفر:} نهى الله المعتكف عن مباشرة النساء حال عكوفه في المسجد.

وذكر \rawi{الطبري} في معنى المباشرة قولين:
\begin{itemize}
    \item أنها الجماع خاصة.
    \item أنها جميع معاني المباشرة من لمس وقبلة وجماع.
\end{itemize}

ثم رجح أن المراد ما كان في معنى الجماع أو قام مقامه مما يوجب الغسل، مستدلاً بما صح عن النبي ﷺ من ترجيل عائشة رضي الله عنها رأسه وهو معتكف، فدل ذلك على أن النهي ليس على عموم كل مباشرة، وإنما على بعض أفرادها.

\begin{fawaid}[السنة تخصص عموم اللفظ القرآني]
من طرق الترجيح عند الطبري أن يرد عموم اللفظ إلى ما خصصته السنة الصحيحة. فلما ثبت أن النبي ﷺ كان يترجّل رأسه وهو معتكف، عُلم أن لفظ ``المباشرة'' في الآية ليس مراداً به جميع ما يصدق عليه الاسم، بل بعضه دون بعض.
\end{fawaid}

\section{تأويل (تلك حدود الله فلا تقربوها)}
فسر \rawi{الطبري} الحدود هنا بالمحارم والأوقات والأحكام التي حدها الله لعباده، ونهاهم عن مقاربتها فضلاً عن الوقوع فيها.

وشدّد الشيخ هنا على معنى عظيم: أن المؤمن لا يهون في نفسه أمر المعاصي والحدود لمجرد وجود الخلاف أو لكون بعض الناس يستخف بها، بل الأصل تعظيم أمر الله ونهيه، والحذر من مسلك المرجئة في تهوين العمل وتحقير أثر المعصية.

\begin{fawaid}[تعظيم حدود الله من حقيقة الإيمان]
من فساد الإيمان أن يهون على العبد أمر العمل الصالح أو شأن المعاصي بدعوى أن هذه ليست من مسائل الإيمان الكبرى. بل تعظيم أمر الله ونهيه، والوقوف عند حدوده، من صلب الإيمان، لا من فضوله.
\end{fawaid}

\separator

\section{تأويل (ولا تأكلوا أموالكم بينكم بالباطل)}
\isnad{قال أبو جعفر:} نهى الله عباده أن يأكل بعضهم مال بعض بالباطل، كالغصب والخيانة والرشوة والدعاوى الفاجرة، وأن يدلي أحدهم بماله إلى الحكام ليتوصل بذلك إلى أخذ حق غيره وهو يعلم.

\subsection{معنى (وتدلوا بها إلى الحكام)}
وقف الشيخ هنا مع البحث الإعرابي في قوله تعالى: \textit{وتدلوا بها إلى الحكام}، وبيّن أن الطبري ذكر فيها وجهين مشهورين:
\begin{itemize}
    \item الجزم على العطف على النهي قبله، فيكون المعنى: لا تأكلوا أموالكم بالباطل ولا تدلوا بها إلى الحكام.
    \item النصب على معنى الصرف، أي: لا تجمعوا بين أكل المال بالباطل والتوصل إلى تثبيته أو اقتطاعه بحكم الحاكم.
\end{itemize}

ومهما اختلف التوجيه الإعرابي، فالمعنى الذي استقر عليه الشرح أن التحيل على أكل المال بالخصومة والرشوة واستعمال صورة القضاء في تمرير الباطل داخل في عموم النهي.

\begin{fawaid}[التحليل النحوي قد يكشف سعة المعنى لا مجرد صورة الإعراب]
لم يكن اهتمام الشيخ هنا بالإعراب لمجرد الصناعة، بل لأن اختلاف التوجيه بين الجزم والنصب يظهر سعة النهي: أهو عن فعلين منفصلين، أم عن الجمع بين الأكل الباطل والتوصل إليه بطريق الحكم؟ وكلاهما يثري فهم الآية.
\end{fawaid}

\separator

\section{تأويل (يسألونك عن الأهلة)}
فسر \rawi{الطبري} الأهلة بأنها المواقيت التي يعرف بها الناس أوقات عباداتهم ومعاملاتهم من حج وصيام وعدد وديون وغيرها.

وبيّن أن القرآن صرف السائلين عن طلب العلة الفلكية أو الصورة الكونية المجردة إلى ما هو أنفع لهم في دينهم ومعاشهم: وهو الوظيفة العملية للأهلة في التوقيت والعبادات والحقوق.

\begin{fawaid}[القرآن يوجه السؤال إلى ما ينفع]
من هدايات هذا الموضع أن الجواب القرآني لا يجيب دائماً بحسب فضول السائل أو صورة سؤاله الأولى، بل قد ينقله إلى المعنى الأهم والأنفع له. فالعبرة هنا ليست بشكل الهلال فقط، بل بما جعله الله له من مصالح الناس.
\end{fawaid}

\section{تأويل (وليس البر بأن تأتوا البيوت من ظهورها)}
ذكر \rawi{الطبري} سبب نزول هذه الآية في قوم من الأنصار كانوا إذا أحرموا لم يدخلوا البيوت من أبوابها، بل من ظهورها، تعظيماً لعادة جاهلية ظنوها من البر.

فبين الله لهم أن البر ليس في هذه الصورة المتوهمة، بل في التقوى، وأن الإتيان من الأبواب هو الموافق للفطرة والشرع.

\begin{fawaid}[إبطال التعبد بالعوائد الجاهلية إذا ألبست ثوب البر]
من أعظم مقاصد هذه الآية هدم التصورات الباطلة التي تلبس لبوس التعبد والديانة من غير حجة من الشرع. فليس كل ما تعظمه النفوس أو توارثته العادات يكون براً وقربة، بل الميزان هو ما شرعه الله.
\end{fawaid}

\separator

\section{تأويل (وقاتلوا في سبيل الله الذين يقاتلونكم)}
جمع \rawi{الطبري} الأقوال في هذه الآية، وذكر أن من أهل التأويل من جعلها في أول الأمر مقيدة بمن يقاتل المسلمين، ثم ذكر من قال بنسخها بآيات أخر أوسع في أمر القتال.

ويظهر من صنيع \rawi{الطبري} هنا عنايته بترتيب النزول ومراحل التشريع في باب الجهاد، لكنه مع ذلك لم يتعجل القول بالنسخ، بل رجح أن النهي عن الاعتداء باقٍ، وأن المراد بالاعتداء هنا: قتل من لا يجوز قتله، كالنساء والصبيان ومن لم ينصب الحرب.

\begin{fawaid}[دعوى النسخ لا تثبت بالاحتمال]
من أقوى ما في هذا الموضع احتجاج الطبري بأن دعوى النسخ في آية محتملة لا تصح بلا حجة بينة. فمجرد وجود قول آخر أو آية أوسع في الباب لا يكفي لإبطال الحكم الأول حتى يقوم الدليل على التعارض الذي لا جمع معه.
\end{fawaid}

\section{تأويل (واقتلوهم حيث ثقفتموهم)}
فسر \rawi{الطبري} ``حيث ثقفتموهم'' بمعنى: حيث وجدتموهم وظفرتم بهم، وذكر أن المراد إخراج من أخرج المؤمنين من ديارهم، والرد عليهم بمثل ما فعلوا.

\section{تأويل (والفتنة أشد من القتل)}
ساق \rawi{الطبري} أقوال السلف في الفتنة، ورجح ما يرجع إلى صرف المؤمن عن دينه أو ردّه إلى الشرك، وأن هذا أعظم من مجرد القتل.

\begin{fawaid}[الفتنة في الدين أعظم من فوات النفس]
من أعظم ما تبينه هذه الآية أن إضلال المؤمن أو صده عن دينه أو حمله على الردة أعظم جرماً من قتله؛ لأن مفسدة الدين أشد من مفسدة ذهاب البدن.
\end{fawaid}

\section{تأويل (ولا تقاتلوهم عند المسجد الحرام حتى يقاتلوكم فيه)}
ذكر \rawi{الطبري} اختلاف القراء في \textit{ولا تقاتلوهم} و\textit{ولا تقتلوهم}، ثم بين أن القتال في الحرم كان له حكم مخصوص في تلك المرحلة، ثم جاء ما وسع الحكم بعد ذلك.

ورجح \rawi{الطبري} قراءة \textit{ولا تقاتلوهم}؛ لأنها أوفق بسياق الإذن بالقتال، وأبعد عن إيهام تعليق جواز الدفع على وقوع القتل منهم أولاً.

\section{تأويل (وقاتلوهم حتى لا تكون فتنة ويكون الدين لله)}
فسر \rawi{الطبري} ذلك بأن القتال إنما شرع حتى تزول الفتنة والشرك، ويخضع الناس لدين الله، وذكر في ذلك الآثار المشهورة والمعاني المؤسسة في هذا الباب.

\section{تأويل (الشهر الحرام بالشهر الحرام والحرمات قصاص)}
بيّن \rawi{الطبري} أن المراد بالشهر الحرام هنا ما يتعلق بحرمة الزمان والمكان والإحرام، وأن الحرمات قصاص، أي: إذا انتهكت من المؤمنين حُق لهم أن يجازوا بمثل ما اعتدي به عليهم في تلك الحرمات.

\section{تأويل (فمن اعتدى عليكم فاعتدوا عليه بمثل ما اعتدى عليكم)}
ذكر \rawi{الطبري} قولين:
\begin{itemize}
    \item أن الآية مكية في مجازاة الأذى بمثله أو الصبر عنه.
    \item وأنها مدنية في سياق القتال بعد فرضه.
\end{itemize}

ثم رجح الثاني؛ لأن السياق كله في القتال، ولأن فرض الجهاد لم يكن قد نزل بمكة بعد. وبيّن أن لفظ ``فاعتدوا'' هنا من باب المشاكلة اللفظية والمجازاة بالمثل، لا من باب الإذن في الظلم.

\begin{fawaid}[السياق الزمني والتشريعي من أقوى وجوه الترجيح]
رجح الطبري مدنية هذه الآية وأنها في باب القتال لأن السياق قبلها وبعدها فيه، ولأن فرض الجهاد إنما كان بعد الهجرة. فمعرفة المرحلة التشريعية والزمن الذي نزلت فيه الآية يعين على رد الأقوال الضعيفة.
\end{fawaid}

\begin{fawaid}[ربط الآية بما قبلها يحسم كثيراً من الخلاف]
لم يرجح الطبري مدنية الآية هنا بخبر مستقل فقط، بل لأن الآيات التي قبلها وبعدها كلها في باب الجهاد والحرمات والقصاص. فالسياق المتصل قد يكون في بعض المواضع أقوى من الرواية المنفردة إذا كان منسجماً معها.
\end{fawaid}

\section{تأويل (واتقوا الله واعلموا أن الله مع المتقين)}
ختم \rawi{الطبري} هذا المقطع بأن الله مع المتقين الذين يقفون عند حدوده، ويؤدون فرائضه، ويتجنبون محارمه.

\separator

وقف الشيخ عند هذا الموضع من سورة البقرة، بعد أن جمع في هذا اللقاء بين أحكام الصيام والاعتكاف، وآداب الوقوف عند الحدود، وأصولاً عظيمة في الجهاد والبر والفتنة والتقوى.
